\chapter{Introduction}
\label{chap:introduction}

\section{Background}

After a major earthquake, the first few hours of response are often shaped by
limited access rather than by lack of need. Survivors may be trapped inside
damaged buildings, under rubble, or in areas surrounded by unstable structures.
At the same time, ground search-and-rescue (SAR) teams may have to work through
blocked roads, collapsed access points, debris, damaged communication
infrastructure, and uncertain safety conditions. Physical rescue therefore
remains a ground-team responsibility, but reaching every detected survivor
immediately is not always possible.

Emergency operations in such conditions are usually coordinated through a
command center that must distribute limited SAR, medical, logistics, and
communication resources across the affected region. The command center does
not only need to know where survivors may be located; it also needs to decide
which reported locations should receive immediate support when responders,
supplies, and access routes are limited.

Unmanned aerial vehicles (UAVs) are useful in this setting because they can
enter hazardous or partially inaccessible areas faster than ground vehicles and
can provide aerial observations to a command center. Recent work has studied
UAV-based disaster reconnaissance, survivor detection, localization, aerial
imagery analysis, and post-earthquake scene understanding
\cite{zaidalkilani2026aiuav,che2025uav5g,sirma2025drespnet}. In the classical
workflow, the UAV mainly acts as a reconnaissance platform: it scans the
affected area, reports possible survivor locations, and helps the command
center build situational awareness.

\begin{figure}[htbp]
    \centering
    \includegraphics[width=0.92\textwidth]{baseline_passive_reconnaissance.png}
    \caption{Baseline passive reconnaissance scenario in which scout UAVs
    report survivor locations to a command center, after which assistance
    depends on the availability and accessibility of ground SAR teams.}
    \label{fig:baseline_passive_reconnaissance}
\end{figure}

Figure~\ref{fig:baseline_passive_reconnaissance} illustrates this baseline
passive reconnaissance process. A scout UAV surveys the earthquake-affected
zone and sends observations to the command center. The observations may include
the suspected survivor location, aerial imagery, and contextual information
about the surrounding damage. After the report is accepted, the command center
depends mainly on ground SAR teams to reach the survivor physically. This
workflow is valuable because it improves situational awareness, but it leaves a
waiting period between detection and rescue.

\section{Motivation}

The key motivation of this thesis comes from the limitation of the passive
workflow. A scan UAV can help identify where a survivor may be located, but it
does not by itself deliver assistance. If ground SAR teams are delayed, the
system may know about a survivor while still being unable to provide immediate
support. This is the operational gap addressed in this work.

The proposed response process therefore separates UAV operations into two
roles. Scout UAVs remain responsible for reconnaissance, sensing, and reporting.
Responder UAVs are introduced as a dedicated delivery layer that can carry
lightweight emergency supplies such as water, food, first-aid materials, or
oxygen-support items to feasible drop-off locations near detected survivors.
This separation is important because reconnaissance and delivery create
different operational requirements. Scout UAVs should preserve flight
endurance, coverage, and sensing continuity, while responder UAVs can be loaded
and routed specifically for relief delivery.

\begin{figure}[htbp]
    \centering
    \includegraphics[width=0.92\textwidth]{proposed_golden_hour_bridge.png}
    \caption{Proposed golden-hour bridge scenario in which responder UAVs
    provide temporary relief delivery between survivor detection and physical
    rescue by ground SAR teams.}
    \label{fig:proposed_golden_hour_bridge}
\end{figure}

Figure~\ref{fig:proposed_golden_hour_bridge} shows the proposed golden-hour
bridge scenario. Compared with the passive reconnaissance baseline, the command
center now has an additional response option. After a survivor-assistance task
is accepted, a responder UAV may be dispatched with suitable relief items. The
responder UAV does not replace the ground SAR team and does not perform
physical rescue. It provides temporary assistance during the critical period
between detection and the arrival of ground responders.

The difficulty is that responder UAV support cannot be planned greedily. A UAV
has limited battery energy, flight time, payload capacity, and onboard relief
items. It must also fly through a feasible three-dimensional route rather than
through buildings, rubble, or no-fly regions. If several survivor-assistance
tasks are waiting at the same time, choosing one task affects whether another
task can still be served within the available resources. In addition, new
survivor reports may arrive while a responder UAV is already in flight. These
conditions motivate a planning model that treats task selection, UAV
assignment, service order, routing, payload use, and energy feasibility as
connected decisions rather than as separate steps.

\section{Problem Statement}

This thesis studies the responder-planning stage that begins after the command
center accepts scout-UAV observations as survivor-assistance tasks. Each task
represents a detected survivor or survivor group requiring a specified set of
lightweight relief items. The exact physical rescue operation is still handled
by ground SAR teams, and the survivor-detection algorithm itself is outside the
scope of this thesis. The focus here is the decision layer that determines how
available responder UAVs should deliver temporary assistance before ground
rescue arrives.

The task-generation process is treated as a command-center decision. A scout
UAV first reports a possible survivor location with supporting observations.
The command center then validates the report, estimates the operational
severity of the case, identifies the required relief items, and records the
corresponding quantities. Only after this process is completed is the report
treated as an eligible survivor-assistance task for responder-UAV planning.

Given a set of survivor-assistance tasks, available responder UAVs, relief-item
requirements, and a feasible 3D flight network, the planning problem is to
determine:

\begin{itemize}
    \item which survivor-assistance tasks should be served;
    \item which responder UAV should serve each selected task;
    \item the order in which the selected tasks should be completed; and
    \item the feasible 3D route followed by each responder UAV.
\end{itemize}

The resulting missions must satisfy:

\begin{itemize}
    \item relief-item requirements;
    \item parcel-count and payload-capacity constraints;
    \item service and travel-time requirements;
    \item battery and reserve-energy constraints; and
    \item safe return to the responder base.
\end{itemize}

Each task is associated with an operational severity and a time-dependent
service value that decreases as assistance is delayed. Therefore, the planning
objective is not simply to minimize flight distance or to serve the nearest
task first. The objective is to maximize the total time-dependent service value
obtained by the responder fleet while respecting the physical and resource
limits of each UAV.

The problem also considers survivor-assistance tasks that arrive after
responder missions have already started. Instead of replanning every active
mission from the beginning, PGBM uses a controlled one-for-one task-replacement
rule. A newly arrived task may replace one uncompleted task in an active
mission only when the replacement is feasible and improves the remaining
mission value. Actions that have already been executed are kept unchanged.

\section{Objectives}

The main objective of this thesis is to develop a resource-aware responder-UAV
planning framework for post-earthquake survivor assistance. The framework is
designed to use limited responder-UAV resources where they are most useful
during the waiting period between survivor detection and ground rescue.

The specific objectives are:

\begin{itemize}
    \item To jointly determine which survivor-assistance tasks are served,
    which responder UAVs serve them, in what order they are served, and which
    feasible 3D routes are used.

    \item To represent time-critical task priority using operational severity
    and a time-dependent service value that decreases with service delay.

    \item To incorporate practical responder limitations including relief-item
    requirements, parcel and payload capacities, route-dependent payload
    evolution, battery consumption, reserve energy, and safe return to base.

    \item To develop a controlled one-for-one task-replacement mechanism for
    considering newly arrived survivor-assistance tasks during active missions.

    \item To evaluate the proposed framework through simulation and compare its
    performance with conventional responder-dispatch strategies.
\end{itemize}

\section{Research Contributions}

The expected contributions of this thesis are summarized as follows:

\begin{itemize}
    \item \textbf{Proactive Golden-Hour Bridge Model (PGBM):}
    A planning framework is introduced for using responder UAVs as a temporary
    assistance bridge between survivor detection and ground rescue.

    \item \textbf{Separation of scout and responder roles:}
    Scout UAVs are used for upstream reconnaissance, while a dedicated
    responder fleet is used for relief delivery. PGBM begins after a survivor
    report has been accepted as an assistance task by the command center.

    \item \textbf{Joint task allocation and 3D routing:}
    Task selection, UAV assignment, service order, and feasible 3D route
    selection are considered together, so the routing and resource effects of
    each assignment are visible during planning.

    \item \textbf{Time-dependent service-value objective:}
    PGBM combines operational severity with service delay and maximizes the
    total value obtained from completed assistance tasks.

    \item \textbf{Resource-aware mission planning:}
    The framework considers relief quantities, parcel and payload capacities,
    payload evolution, battery-energy consumption, reserve energy, and safe
    return to base.

    \item \textbf{Controlled dynamic recourse:}
    A one-for-one task-replacement mechanism allows one newly arrived task to
    replace one uncompleted task when the replacement is feasible and improves
    the remaining mission value.
\end{itemize}

This scope is intentionally limited. PGBM does not propose a new
survivor-detection, environmental-mapping, or physical-rescue algorithm.
Instead, it operates as a downstream decision layer that uses accepted survivor
information and an externally supplied feasible 3D free-space representation.

\section{Research Challenges}

The main challenges addressed in this research include:

\begin{itemize}
    \item \textbf{Joint decision making:}
    Task assignment and routing are strongly coupled because route length,
    completion time, payload, and battery usage influence which assignments
    are feasible.

    \item \textbf{Time-sensitive prioritization:}
    Higher-priority tasks should generally receive earlier assistance, but the
    planner must consider the effect of each decision on the total mission
    value.

    \item \textbf{Three-dimensional route feasibility:}
    A detected survivor may be located inside rubble or a damaged structure,
    so the responder must use a nearby feasible drop-off waypoint and travel
    only through permitted 3D free space.

    \item \textbf{Changing payload and energy use:}
    Relief payload decreases after deliveries, causing the UAV's carried mass
    and payload-dependent energy consumption to change throughout the route.

    \item \textbf{Dynamic task arrival:}
    New survivor-assistance tasks may appear after dispatch, requiring mission
    adaptation without changing already completed actions.

    \item \textbf{Limited onboard supplies:}
    An active responder can serve a newly arrived task only when its current
    onboard inventory is sufficient for the modified remaining mission.
\end{itemize}

\section{Organization of the Thesis}

The remainder of this thesis is organized as follows:

\begin{itemize}

\item \textbf{
\hyperref[chapter_related_work]
{Chapter~\ref*{chapter_related_work} -- Related Work:}
}
    Reviews existing research on UAV-assisted disaster response,
    survivor detection and localization, multi-UAV search, 3D coverage,
    resource-aware mission planning, and adaptive UAV coordination.

    \item \textbf{\hyperref[chapter_proposedMethod]
    {Chapter~\ref*{chapter_proposedMethod} -- Proposed Methodologies:}}
    Presents the PGBM system model and mathematical problem formulation,
    including the time-dependent service-value objective, joint task allocation
    and 3D routing, payload and battery constraints, and the controlled
    one-for-one task-replacement mechanism. It also describes the overall
    operational workflow, responder-mission planning, mission execution,
    state updates, and dynamic recourse procedure.

    \item \textbf{\hyperref[chapter_implementation]
    {Chapter~\ref*{chapter_implementation} -- Implementation:}}
    Describes the simulation environment, system parameters, optimization
    implementation, and experimental setup used to evaluate PGBM.

    \item \textbf{\hyperref[chapter_experimentalResults]
    {Chapter~\ref*{chapter_experimentalResults} -- Experimental Results:}}
    Evaluates PGBM under different task-arrival, responder, environmental,
    and resource conditions and compares its performance with selected
    responder-dispatch strategies.

    \item \textbf{\hyperref[chapter_conclusion]
    {Chapter~\ref*{chapter_conclusion} -- Conclusion:}}
    Summarizes the main findings, discusses the limitations of the proposed
    framework, and presents possible directions for future research.

\end{itemize}
